% ============================================================================
% TROUBLESHOOTING AND KNOWN GOTCHAS
% ============================================================================

\section{Troubleshooting and Known Gotchas}
\label{section:troubleshooting}

This section is the lessons-learned dump. Every entry below cost
real bench-debug time; if you hit a similar symptom, read the
corresponding entry first before reaching for an oscilloscope.

\subsection{USB / CDC enumeration}

\subsubsection{ttyACM* indices change after re-flash}

\textbf{Symptom:} after \texttt{rfp-cli} reflash the gateway
fails to open \texttt{/dev/ttyACM1} (or whichever index it had
before).

\textbf{Why:} the firmware re-init cycles the USBb peripheral, the
host sees a disconnect + re-enumerate, and Linux assigns the next
free index. Hot-plugging another CDC device between the cycles
shifts the numbers further.

\textbf{Fix:} install the udev symlinks:
\texttt{sudo bash scripts/setup-pi5-host.sh} (which deploys
\texttt{scripts/71-star-symlinks.rules}). The gateway can then
open \texttt{/dev/star-rx72n-proto / -decoded / -log} and not
care what \texttt{ttyACM*} number it ends up as.

\subsubsection{No CY7C65213 UART bridge while E2 Lite is plugged in}

\textbf{Symptom:} \texttt{ls /dev/cu.usbmodem* /dev/cu.usbserial-*}
shows no Cypress-bridge serial device while the E2 Lite is connected.

\textbf{Why:} bench harness shares USB power gating between the
E2 Lite and the Cypress bridge. This is hardware wiring, not a
fault.

\textbf{Fix:} disconnect the E2 Lite physically after flashing
before running anything that opens the SCI9 fallback path. The
USB0 path on \texttt{/dev/ttyACM*} is independent and does \emph{not}
suffer from this constraint.

\subsubsection{TOM enumerates fine on Pi 5 despite HUM warnings}

\textbf{Symptom:} The Renesas HW manual says PLLSRCSEL must be 0
when using USB; TOM violates this (PLLSRCSEL=1, HOCO source) and
HOCO accuracy is $\pm$2.44\,\% vs USB FS spec $\pm$0.25\,\%.

\textbf{Reality:} The Pi 5 xHCI controller accepts the HOCO-driven
clock without issue. Bench-verified 2026-04-24 on commit
\texttt{341bd09ce}: TOM enumerates the 3-port CDC composite
cleanly, idVendor=045b idProduct=0235, no EPIPE / EPROTO.

\textbf{Caveat:} stricter xHCI variants (some macOS revisions,
some Windows HCDs, embedded controller hubs) MAY reject this clock.
For ship-to-Pi deployments, TOM USB is fine. For field deployments
on unknown hosts, the 24\,MHz crystal-clocked PROD configuration
is the safer choice.

\subsubsection{USB ISR not firing despite ICU/IER configured}

\textbf{Symptom:} \texttt{g\_usb\_isr\_entry\_count == 0} even
though INTSTS0 / BRDY / BEMP latch correctly inside the USB peripheral.
Vector 144 (USBI0) configured cleanly; nothing fires.

\textbf{Why:} two HUM-mandated steps were missing.

\begin{enumerate}
  \item \textbf{SLIPRCR.WPRC was never written.} HUM 15.7.7 step (5)
        makes \texttt{SLIPRCR.WPRC = 1} mandatory after assigning a
        software-configurable interrupt. Without it the ICU silently
        drops USBI0.
  \item \textbf{SYSCFG.USBE was set BEFORE SYSCFG.SCKE.} HUM
        40.3.1.1 requires SCKE first; reverse order silently no-ops
        half the SYSCFG follow-on writes because the USB module
        clock is gated until SCKE=1 latches.
\end{enumerate}

\textbf{Fix:} both fixes shipped on \texttt{feat/usb0} commit
\texttt{f23b15eb5}; they are now in main. If you see this symptom
again, check that those two specific writes are present and in
the right order in your bring-up code.

\subsection{Bench-rig pitfalls}

\subsubsection{AD2 Single-mode buffer is 4096 samples, period}

\textbf{Symptom:} an AD2 capture longer than ~4 ms appears to show
the MUT going dead. You set
\texttt{FDwfDigitalInBufferSizeSet(8192)} and run a 20-second
sweep; only the first 4\,ms are useful.

\textbf{Why:} \texttt{FDwfDigitalInBufferSizeSet(n)} silently
clamps to 4096 samples no matter what you ask for. With
\texttt{SampleFormat=8} (one byte per word), 4096 samples at
1\,MHz = ~4\,ms (sometimes ~8\,ms with byte-format packing). The
AD2 isn't recording past that point, and your firmware looks
mute.

\textbf{Fix:} pick the sample rate to match the desired window:
\texttt{sample\_hz = buffer\_size / desired\_seconds}. For a 20\,s
debug-trace, that's 200\,Hz. For high-frequency signals, switch to
Record mode (\texttt{FDwfDigitalInAcquisitionModeSet mode=2}) which
streams continuously to the host.

\subsubsection{sci9\_debug\_puts hangs from a ThreadX task}

\textbf{Symptom:} debug logging via \texttt{sci9\_debug\_puts}
freezes the calling task; everything else continues but that task
is wedged.

\textbf{Why:} the helper busy-polls TDRE before writing. SCI9 runs
on PCLKA (120\,MHz) but its BRR was being computed against PCLKB
(60\,MHz) -- so the configured baud was halved and TDRE never
came back true at the expected rate within the busy-poll budget.

\textbf{Fix until BRR/PCLK is corrected:} skip \texttt{sci9}
debug from any ThreadX task; use the rx\_log path or USB CDC log
port 3 instead.

\subsection{Build / toolchain}

\subsubsection{Devcontainer rebuild loop}

\textbf{Symptom:} VS Code rebuilds the devcontainer image every
session; takes minutes.

\textbf{Why:} the GHCR cache is missing or stale for the user.

\textbf{Fix:} pull the cached image manually:
\texttt{docker pull ghcr.io/locked-inc/star/gnurx:latest}. The
\texttt{firmware-toolchain-image.yml} workflow keeps it warm.

\subsubsection{Local rx-elf-gcc fails; should I use the devcontainer?}

\textbf{Memory rule:} use the devcontainer only if the local
\texttt{rx-elf-gcc} actually fails. If it works, build directly --
it's faster. There is no \texttt{devcontainer-exec.sh} script and
you should not create one.

\subsection{Coding-policy reminders}

\subsubsection{No NOLINT for readability-* clang-tidy checks}

\textbf{Rule:} \texttt{readability-function-size},
\texttt{readability-function-cognitive-complexity},
\texttt{readability-magic-numbers}, etc. must be REFACTORED, never
suppressed with \texttt{NOLINT}.

\textbf{Why:} the user's standing rule. Suppressions accumulate
silently and the codebase loses the structural feedback the lint
gave for free.

\textbf{Exception:} path-sensitive checks like
\texttt{clang-analyzer-optin.core.EnumCastOutOfRange} where there
is genuinely no source-level workaround -- a NOLINT is acceptable
there. \texttt{readability-*} is not in that exception class.

\subsubsection{Trust headers over comments}

\textbf{Rule:} when a stripped-down test file's comment and the
\texttt{libs/rx\_hal/inc/*.h} header disagree, the header wins. The
clock tree comment in \texttt{encoder\_test/} (or anywhere)
claiming PCLKB = 120\,MHz is wrong; \texttt{rx72n\_clock.h} says
60\,MHz and that is the truth.

\subsubsection{Trust known-good bench tests over fresh code}

\textbf{Rule:} if \texttt{motor\_spin\_test/}, \texttt{encoder\_test/},
or \texttt{imu\_test/} already initialize a peripheral successfully
on the bench, diff your new code against theirs. Do not freelance
an init sequence when a proven one exists. CLAUDE.md elevates this
to a P0 rule because guessing-and-flashing has cost the project
real days.

\subsection{Common red herrings}

\subsubsection{"It looks like the MCU halts 8\,ms after rfp-cli disconnect"}

It does not. The AD2 buffer simply ran out (see 4096-sample issue
above). Re-sample at \texttt{buffer\_size / desired\_seconds} and
the trace will run for the full window.

\subsubsection{"USB enumeration is broken because BRR is wrong"}

Not necessarily. Before blaming firmware, verify the host side:

\begin{verbatim}
stty -F /dev/ttyACM1 -a
\end{verbatim}

Old stty state from a previous session can sit on a tty and look
exactly like a baud-rate mismatch. CLAUDE.md's "Embedded Gotchas"
list calls this out for the SCI9 path; the same lesson applies to
USB CDC re-enumerations.

\subsubsection{"The peripheral isn't responding to register writes"}

Most likely it is in module-stop. The MSTPCRA / MSTPCRB / MSTPCRC /
MSTPCRD bit for that peripheral has to be cleared (\texttt{0}) before
any of its registers are addressable. Silent failure mode.

Other common silent failures:

\begin{itemize}
  \item \textbf{PFS locked}: writes to \texttt{P\{n\}PFS} require
        \texttt{PWPR.B0WI=0} then \texttt{PWPR.PFSWE=1}. Missing
        unlock = pin-mux never changes = no signal at the pad.
  \item \textbf{MPC PSEL wrong}: PSEL value differs per pin. Check
        the MPC chapter's pin-function table, not a sibling pin.
  \item \textbf{Wrong PCLK}: SCI uses PCLKB; GPTW uses PCLKA;
        ADC12 uses PCLKD. Mixing silently miscomputes the divider.
  \item \textbf{FPU context save not enabled}: a task using floats
        without \texttt{TX\_THREAD\_EXTENSION\_2} silently corrupts
        another task's FPU state on preemption.
  \item \textbf{Encoder TMDR=0 instead of phase-counting mode}:
        gives a free-running timer, not an encoder.
\end{itemize}

CLAUDE.md's "Embedded Gotchas to Pattern-Match Instantly" list is
the expanded version; this section is the curated subset that has
actually bitten this project.

\subsection{Where to look in code}

\begin{center}
\begin{tabular}{|p{5.5cm}|p{8.5cm}|}
\hline
\textbf{Topic}                                & \textbf{Source}                                                                          \\
\hline
udev symlink rules                            & \texttt{scripts/71-star-symlinks.rules}                                                  \\
USB autosuspend disable                       & \texttt{scripts/setup-pi5-host.sh} (installs \texttt{70-star-usb-no-autosuspend.rules}) \\
USB ISR routing fix                           & \texttt{star-rx72n-firmware/libs/rx\_usb/src/rx\_usb\_hw.c}                              \\
SCI9 debug helper                             & \texttt{star-rx72n-firmware/encoder\_test/sci9.c} (and similar in sibling bench tests)   \\
AD2 capture harness                           & \texttt{star-rx72n-firmware/gpio\_test/host/}                                            \\
Embedded gotchas (canonical)                  & \texttt{CLAUDE.md} "Engineering Discipline" section                                       \\
\hline
\end{tabular}
\end{center}
